\section{Mechanized core, and what the finite result is an instance of}

Everything above this line is an enumeration over an authored state space. That is a true statement about 3,072 states and thirteen donor families, and it is not a theorem. Three machine-checked artifacts under \idt{formal/mechanized/} now say which parts of it follow from a general semantics, and what each part costs.

\textbf{The calculus.} Ten composition and revocation theorems --- scope non-amplification, domain confinement, non-laundering, non-compensatory obligations, defeater monotonicity, epoch isolation, the one-step delegation lemma, a bridge from conversion to reachability, and a characterisation of cycles --- are discharged by Z3 over uninterpreted domain, object and issuer sorts, so none of them is a statement about thirteen donors or six named domains. Two limits travel with them and are recorded in the artifact rather than in a footnote: reachability is \emph{axiomatised} as a reflexive-transitive closure carrying a well-founded rank, because transitive closure is not first-order definable, and the chain theorem's induction schema is the single hand step in the development. The rank is not a formality --- an earlier axiomatisation without it admitted circular reachability, in which a sink domain reached its own source, and a differential against the committed checker is what caught it.

\textbf{The state model as an instance.} Each of the seven arguments the X4 terminal takes is assigned to the calculus conjunct it means: type-coordinate agreement or a registered protected coercion is the calculus's reachability, narrowing is scope containment, and the blocker is at once a hard obligation and a defeater. The calculus's own rule under that assignment reproduces X4 exactly on all 3,072 distinct states, with all four terminals reached.

\textbf{The chain compositions as instances.} A donor family is interpreted as an element of an uninterpreted sort whose \emph{domain is its scientific type profile}; a protected coercion is a registered conversion and narrowing is scope containment. Under eight frame conditions stated as axioms --- each shown load-bearing by exhibiting a countermodel to a theorem it carries --- seven theorems are discharged, and donor-level chain confinement is expanded and discharged at every chain length up to six. The bound is a measured limit rather than a chosen one: lengths seven to twelve were tried and the solver returned \texttt{unknown} on some runs and a proof in seconds on others, so the ladder stops where it is a reliable corroboration. Nothing about the theorem changes at length seven --- the induction schema is what carries every length, and it is the single hand step in the development. A this paper chain hop \emph{is} the calculus's delegation, which is itself one of the discharged theorems, so Theorem V3.7 is an instance of a statement that mentions neither thirteen nor two.

\textbf{What that costs Theorem V3.7.} The derivation lands, and what it lands on is small. We state it here rather than leave a reader to assume otherwise.

\begin{itemize}
  \item The 169 is one composition counted 169 times. The committed enumeration writes its chain claim as a doubly nested loop over the thirteen families whose body mentions neither loop variable, and under the interpretation all 169 ordered pairs compose to a single distinct state. The donor axis is a replication factor here exactly as it is in the 39,936 evaluations, and here it is squared.
  \item Nothing heterogeneous is exercised. The state the chain claim evaluates has all five type coordinates agreeing, which by the interpretation puts both donors in one domain; the hop is then a reach by reflexivity and no conversion is ever consulted. Read the thirteen families as type-distinct, so that a cross-family hop crosses a type boundary unless a protected bridge is registered, and the same committed rule returns \textbf{13} compositions rather than 169. Registering every cross-family bridge returns 169 again. The published number therefore records the reading under which no two of the thirteen donor families differ in scientific type --- which is the reading under which the chain is not heterogeneous.
  \item Neither published count tests the interpretation. All eight wrong composition operators tried reproduce 169 successes exactly, and six of the eight also reproduce the 169 widening countermodels. What separates them is an exhaustive identity checked through the committed rule on 36,864 representative pairs, which stand in exactly for all 9,437,184 pairs of X4 states: a composed chain discharges when and only when both of its hops do. Every one of the eight wrong operators breaks it, and the shipped rule satisfies it with no exception.
\end{itemize}

None of the above has been checked by anyone outside the lane that wrote it. Independent formal review and independent systems reproduction remain open, and no empirical claim is made or supported here.
